\section{Exemplos Integrados}

Esta seção apresenta exercícios resolvidos do livro \textit{Estatística Fácil}, desenvolvidos passo a passo.

\subsection{Distribuição Binomial}

Um teste de múltipla escolha contém 6 questões, cada uma com 4 alternativas e apenas uma correta ($p = 0{,}25$). Se um candidato responder a todas as questões de forma aleatória (chute), determine a probabilidade de ele acertar exatamente 2 questões.

\subsubsection*{Resolução Manual}

A probabilidade de obter $k$ sucessos em $n$ ensaios independentes em uma Distribuição Binomial é dada por:
\[
P(X = k) = \binom{n}{k} \cdot p^k \cdot (1-p)^{n-k}
\]

Identificando os parâmetros do problema:
\begin{itemize}
    \item Número de questões ($n$) = $6$
    \item Número de acertos desejado ($k$) = $2$
    \item Probabilidade de acerto ($p$) = $0{,}25$
    \item Probabilidade de erro ($1-p$) = $1 - 0{,}25 = 0{,}75$
\end{itemize}

1. Cálculo da combinação simples $\binom{6}{2}$:
\[
\binom{6}{2} = \frac{6!}{2!(6-2)!} = \frac{6 \times 5}{2 \times 1} = 15
\]

2. Substituição dos valores na fórmula binomial:
\[
P(X = 2) = 15 \cdot (0{,}25)^2 \cdot (0{,}75)^{6-2}
\]
\[
P(X = 2) = 15 \cdot 0{,}0625 \cdot (0{,}75)^4 = 15 \cdot 0{,}0625 \cdot 0{,}31640625 \approx 0{,}2966
\]

\subsubsection*{Resolução na Calculadora}

O cálculo direto de probabilidades binomiais é realizado no modo de cálculo geral \textbf{COMP} (\texttt{MODE 1}), utilizando a tecla dedicada \textbf{nCr} para a combinação e a tecla \textbf{$\wedge$} para as potências.

\begin{table}[H]
    \centering
    \begin{tabular}{p{4.5cm} p{5.5cm} p{3cm}}
        \hline
        \textbf{Etapa} & \textbf{Entrada} & \textbf{Resultado} \\ \hline
        Ativar modo COMP & \texttt{MODE 1} & \textit{Modo Geral} \\
        Cálculo da combinação $\binom{6}{2}$ & \texttt{6 nCr 2 =} & \texttt{15} \\
        Cálculo direto de $P(X=2)$ & \texttt{(6 nCr 2) * (0.25 \^{} 2) * (0.75 \^{} ( 6 - 2 )) =} & \texttt{0.296630859} \\ \hline
    \end{tabular}
\end{table}

A probabilidade de o candidato acertar exatamente duas questões chutando é de aproximadamente $29{,}66\%$.

\subsection{Variância e Desvio Padrão Amostral}

Considerando a amostra das estaturas (em centímetros) de 5 estudantes: $\{160, 165, 170, 172, 178\}$, determine a variância amostral ($s^2$) e o desvio padrão amostral ($s$).

\subsubsection*{Resolução Manual}

1. Cálculo da média aritmética ($\bar{x}$):
\[
\bar{x} = \frac{\sum x_i}{n} = \frac{160 + 165 + 170 + 172 + 178}{5} = \frac{845}{5} = 169 \text{ cm}
\]

2. Cálculo dos desvios ao quadrado e da soma dos quadrados ($\sum (x_i - \bar{x})^2$):
\begin{align*}
(160 - 169)^2 &= (-9)^2 = 81 \\
(165 - 169)^2 &= (-4)^2 = 16 \\
(170 - 169)^2 &= (1)^2 = 1 \\
(172 - 169)^2 &= (3)^2 = 9 \\
(178 - 169)^2 &= (9)^2 = 81
\end{align*}
\[
\sum_{i=1}^{5} (x_i - \bar{x})^2 = 81 + 16 + 1 + 9 + 81 = 188
\]

3. Cálculo da variância amostral ($s^2$):
\[
s^2 = \frac{\sum (x_i - \bar{x})^2}{n - 1} = \frac{188}{5 - 1} = \frac{188}{4} = 47 \text{ cm}^2
\]

4. Cálculo do desvio padrão amostral ($s$):
\[
s = \sqrt{s^2} = \sqrt{47} \approx 6{,}8556 \text{ cm}
\]

\subsubsection*{Resolução na Calculadora}

Utiliza-se o modo estatístico de uma variável \textbf{SD} (\texttt{MODE 2}).

\begin{table}[H]
    \centering
    \begin{tabular}{p{4.5cm} p{5.5cm} p{3cm}}
        \hline
        \textbf{Etapa} & \textbf{Entrada} & \textbf{Resultado} \\ \hline
        Ativar modo SD & \texttt{MODE 2} & \texttt{SD} \\
        Limpar memória estatística & \texttt{SHIFT CLR 1 =} & \texttt{Stat-Clear} \\
        Inserir dados & \texttt{160 DT 165 DT 170 DT 172 DT 178 DT} & \texttt{n=5} \\
        Recuperar Média ($\bar{x}$) & \texttt{SHIFT 2 1 =} & \texttt{169} \\
        Desvio Padrão Amostral ($s$) & \texttt{SHIFT 2 3 =} & \texttt{6.855654601} \\
        Variância Amostral ($s^2$) & \texttt{Ans $x^2$ =} & \texttt{47} \\ \hline
    \end{tabular}
\end{table}

\subsection{Média e Desvio Padrão com Tabela de Frequências}

A tabela a seguir apresenta a distribuição do número de acertos de 20 candidatos em um teste de conhecimento:

\begin{table}[H]
    \centering
    \begin{tabular}{cc}
        \hline
        \textbf{Número de Acertos ($x_i$)} & \textbf{Frequência ($f_i$)} \\ \hline
        2 & 2 \\
        4 & 5 \\
        6 & 8 \\
        8 & 3 \\
        10 & 2 \\ \hline
    \end{tabular}
    \label{tab:frequencias_exemplo}
\end{table}

Com base na tabela de frequências apresentada, calcule a média aritmética ($\bar{x}$) e o desvio padrão da distribuição.

\subsubsection*{Resolução Manual}

Como os dados representam o total de candidatos do teste, aplica-se a fórmula populacional para o desvio padrão.

1. Tabela auxiliar para os somatórios ($\sum f_i$, $\sum f_i x_i$ e $\sum f_i x_i^2$):

\begin{table}[H]
    \centering
    \begin{tabular}{ccccc}
        \hline
        $x_i$ & $f_i$ & $f_i x_i$ & $x_i^2$ & $f_i x_i^2$ \\ \hline
        2  & 2  & 4   & 4   & 8   \\
        4  & 5  & 20  & 16  & 80  \\
        6  & 8  & 48  & 36  & 288 \\
        8  & 3  & 24  & 64  & 192 \\
        10 & 2  & 20  & 100 & 200 \\ \hline
        \textbf{Total} & \textbf{20} & \textbf{116} & -- & \textbf{768} \\ \hline
    \end{tabular}
\end{table}

2. Média aritmética ($\bar{x}$):
\[
\bar{x} = \frac{\sum f_i x_i}{\sum f_i} = \frac{116}{20} = 5{,}8
\]

3. Desvio padrão populacional:
\[
s = \sqrt{\frac{\sum f_i x_i^2}{n} - \bar{x}^2} = \sqrt{\frac{768}{20} - (5{,}8)^2}
\]
\[
s= \sqrt{38{,}4 - 33{,}64} = \sqrt{4{,}76} \approx 2{,}1817
\]

\subsubsection*{Resolução na Calculadora}

No modo \textbf{SD} (\texttt{MODE 2}), os dados agrupados por frequência são digitados utilizando o separador \texttt{;} (obtido por \texttt{SHIFT ,}).

\begin{table}[H]
    \centering
    \begin{tabular}{p{4.5cm} p{5.5cm} p{3cm}}
        \hline
        \textbf{Etapa} & \textbf{Entrada} & \textbf{Resultado} \\ \hline
        Ativar modo SD & \texttt{MODE 2} & \texttt{SD} \\
        Limpar dados estatísticos & \texttt{SHIFT CLR 1 =} & \texttt{Stat-Clear} \\
        Inserir $x_1=2, f_1=2$ & \texttt{2 ; 2 DT} & \texttt{n=2} \\
        Inserir $x_2=4, f_2=5$ & \texttt{4 ; 5 DT} & \texttt{n=7} \\
        Inserir $x_3=6, f_3=8$ & \texttt{6 ; 8 DT} & \texttt{n=15} \\
        Inserir $x_4=8, f_4=3$ & \texttt{8 ; 3 DT} & \texttt{n=18} \\
        Inserir $x_5=10, f_5=2$ & \texttt{10 ; 2 DT} & \texttt{n=20} \\
        Recuperar Média ($\bar{x}$) & \texttt{SHIFT 2 1 =} & \texttt{5.8} \\
        Desvio Padrão Pop. ($\sigma$) & \texttt{SHIFT 2 2 =} & \texttt{2.181742423} \\ \hline
    \end{tabular}
\end{table}

\subsection{Regressão Linear Simples}

Para estudar a relação entre os gastos em publicidade ($x$, em milhares de reais) e o volume de vendas ($y$, em milhares de reais) de uma empresa, foram coletados os seguintes pares de dados referentes a 5 meses:

\begin{table}[H]
    \centering
    \begin{tabular}{cc}
        \hline
        \textbf{Publicidade ($x$)} & \textbf{Vendas ($y$)} \\ \hline
        2 & 4 \\
        3 & 5 \\
        5 & 7 \\
        6 & 8 \\
        8 & 9 \\ \hline
    \end{tabular}
    \label{tab:regressao_exemplo}
\end{table}

Obtenha a equação da reta de regressão linear simples ($\hat{y} = A + Bx$) e o coeficiente de correlação linear ($r$).

\subsubsection*{Resolução Manual}

1. Tabela auxiliar para os somatórios de regressão ($n = 5$):

\begin{table}[H]
    \centering
    \begin{tabular}{ccccc}
        \hline
        $x_i$ & $y_i$ & $x_i^2$ & $y_i^2$ & $x_i y_i$ \\ \hline
        2 & 4 & 4  & 16 & 8   \\
        3 & 5 & 9  & 25 & 15  \\
        5 & 7 & 25 & 49 & 35  \\
        6 & 8 & 36 & 64 & 48  \\
        8 & 9 & 64 & 81 & 72  \\ \hline
        \textbf{24} & \textbf{33} & \textbf{138} & \textbf{235} & \textbf{178} \\ \hline
    \end{tabular}
\end{table}

2. Cálculo da inclinação (coeficiente angular $B$):
\[
B = \frac{n \sum x_i y_i - \sum x_i \sum y_i}{n \sum x_i^2 - (\sum x_i)^2} = \frac{5(178) - (24)(33)}{5(138) - (24)^2}
\]
\[
B = \frac{890 - 792}{690 - 576} = \frac{98}{114} \approx 0{,}8596
\]

3. Cálculo do intercepto (coeficiente linear $A$):
\[
\bar{x} = \frac{24}{5} = 4{,}8 \quad \text{e} \quad \bar{y} = \frac{33}{5} = 6{,}6
\]
\[
A = \bar{y} - B \bar{x} = 6{,}6 - (0{,}8596 \cdot 4{,}8) = 6{,}6 - 4{,}1261 = 2{,}4739
\]

A equação da reta de regressão é:
\[
\hat{y} = 2{,}4739 + 0{,}8596x
\]

4. Cálculo do coeficiente de correlação linear ($r$):
\[
r = \frac{n \sum x_i y_i - \sum x_i \sum y_i}{\sqrt{\left[ n \sum x_i^2 - (\sum x_i)^2 \right] \left[ n \sum y_i^2 - (\sum y_i)^2 \right]}}
\]
\[
r = \frac{98}{\sqrt{114 \cdot \left[ 5(235) - (33)^2 \right]}} = \frac{98}{\sqrt{114 \cdot (1175 - 1089)}} = \frac{98}{\sqrt{114 \cdot 86}}
\]
\[
r = \frac{98}{\sqrt{9804}} = \frac{98}{99{,}0151} \approx 0{,}9897
\]

O valor $r \approx 0{,}9897$ indica uma correlação linear positiva muito forte entre o investimento em publicidade e o volume de vendas.

\subsubsection*{Resolução na Calculadora}

A análise de regressão linear exige a ativação do modo \textbf{REG LIN} (\texttt{MODE 3 1}). Os pares de dados $(x, y)$ são informados separando as variáveis por vírgula \texttt{,} e confirmados com \texttt{DT}.

\begin{table}[H]
    \centering
    \begin{tabular}{p{4.5cm} p{5.5cm} p{3cm}}
        \hline
        \textbf{Etapa} & \textbf{Entrada} & \textbf{Resultado} \\ \hline
        Ativar Modo REG LIN & \texttt{MODE 3 1} & \texttt{REG} \\
        Limpar memória estatística & \texttt{SHIFT CLR 1 =} & \texttt{Stat-Clear} \\
        Inserir par (2, 4) & \texttt{2 , 4 DT} & \texttt{n=1} \\
        Inserir par (3, 5) & \texttt{3 , 5 DT} & \texttt{n=2} \\
        Inserir par (5, 7) & \texttt{5 , 7 DT} & \texttt{n=3} \\
        Inserir par (6, 8) & \texttt{6 , 8 DT} & \texttt{n=4} \\
        Inserir par (8, 9) & \texttt{8 , 9 DT} & \texttt{n=5} \\
        Coeficiente $A$ (intercepto) & \texttt{SHIFT 2 1 =} & \texttt{2.473684211} \\
        Coeficiente $B$ (inclinação) & \texttt{SHIFT 2 2 =} & \texttt{0.859649122} \\
        Coeficiente de correlação ($r$) & \texttt{SHIFT 2 3 =} & \texttt{0.989749437} \\ \hline
    \end{tabular}
\end{table}